\documentclass[11pt]{article}
\usepackage[a4paper,margin=1in]{geometry}
\usepackage{graphicx}
\usepackage{svg}
\usepackage{hyperref}
\usepackage{enumitem}
\usepackage{array}
\usepackage{longtable}
\usepackage{xcolor}
\usepackage{caption}

\hypersetup{
  colorlinks=true,
  linkcolor=black,
  urlcolor=blue
}

\title{\textbf{System Design Document}\\AntriView AI}
\author{Rishiwant Kumar Maurya}
\date{April 2026}

\begin{document}
\maketitle
\tableofcontents
\newpage

% ------------------------------------------------------------
\section{Introduction}

\subsection{Overview of the System}
AntriView AI is a premium, UI-first mock interview preparation platform built as a single-page web application.
It provides a guided interview flow (setup \(\rightarrow\) timed questions \(\rightarrow\) session report), a progress dashboard, and supporting tools such as persona selection, resume analysis (demo), and account settings.

\subsection{Problem Statement}
Interview preparation is often fragmented across multiple tools (question banks, timers, notes, and feedback templates).
Candidates struggle to maintain consistency, measure improvement, and practice in conditions that resemble real interviews.

\subsection{Objectives of the System}
\begin{itemize}[leftmargin=1.2em]
  \item Provide an end-to-end mock interview experience: configure a session, practice, and review outcomes.
  \item Track progress over time (sessions, skill scores, and history) in a user-centric dashboard.
  \item Maintain a clean architecture structure (domain/application/infrastructure/presentation) suitable for academic/portfolio use.
  \item Keep the system lightweight and easy to run locally using a modern frontend toolchain.
\end{itemize}

% ------------------------------------------------------------
\section{System Architecture}

\subsection{High-level Architecture}
The system is a client-side application structured into layered modules:
\begin{itemize}[leftmargin=1.2em]
  \item \textbf{Presentation Layer (UI)}: React components in \texttt{src/components/} render pages and user interactions.
  \item \textbf{Application Layer (Use Cases)}: business logic in \texttt{src/application/useCases/} (auth, sessions, profile).
  \item \textbf{Domain Layer (Entities/Ports)}: core types and interfaces in \texttt{src/domain/}.
  \item \textbf{Infrastructure Layer (Repositories/Storage)}: localStorage/sessionStorage adapters in \texttt{src/infrastructure/}.
  \item \textbf{Composition / DI Container}: \texttt{src/app/services.ts} wires dependencies; \texttt{ServicesProvider} exposes them via React Context.
\end{itemize}

\subsection{Frontend, Backend, Database, External Services}
\begin{itemize}[leftmargin=1.2em]
  \item \textbf{Frontend}: React 19 + TypeScript, bundled with Vite; UI uses CSS variables + glassmorphism styling and Lucide icons.
  \item \textbf{Backend}: Not present in the current version. ``AI'' behavior is simulated in the UI; \texttt{src/lib/gemini.ts} is a placeholder for future secure integration.
  \item \textbf{Database}: Browser storage.
  \begin{itemize}[leftmargin=1.2em]
    \item \textbf{User DB}: JSON stored in localStorage key \texttt{antriview\_user\_db}.
    \item \textbf{Session}: current user email + view state stored in sessionStorage keys \texttt{antriview\_current\_user}, \texttt{antriview\_view}, \texttt{antriview\_dash\_view}.
  \end{itemize}
  \item \textbf{External Services / Browser APIs}:
  \begin{itemize}[leftmargin=1.2em]
    \item MediaDevices API for webcam preview (\texttt{getUserMedia}) in the interview room.
    \item Web Speech API capability exists as a utility (\texttt{src/lib/speech.ts}); the interview transcript is currently simulated.
  \end{itemize}
\end{itemize}

\subsection{Component Interaction}
\begin{itemize}[leftmargin=1.2em]
  \item UI components call \textbf{use cases} via \texttt{useServices()}.
  \item Use cases operate on \textbf{ports} (\texttt{UserRepository}, \texttt{SessionRepository}).
  \item Infrastructure repositories persist data to browser storage (localStorage/sessionStorage).
  \item The \texttt{Dashboard} orchestrates navigation between internal views (overview, setup, interview, report, etc.).
\end{itemize}

% ------------------------------------------------------------
\section{Design Principles Used}

\subsection{SOLID Principles}
\begin{itemize}[leftmargin=1.2em]
  \item \textbf{Single Responsibility}: Use cases focus on one domain action (e.g., \texttt{createAuthUseCases}, \texttt{createSessionUseCases}, \texttt{createUserProfileUseCases}).
  \item \textbf{Open/Closed}: Repositories are behind interfaces (ports). Alternate implementations (e.g., server-backed repositories) can be introduced without rewriting use cases.
  \item \textbf{Liskov Substitution}: Any implementation of \texttt{UserRepository} can replace \texttt{LocalStorageUserRepository}.
  \item \textbf{Interface Segregation}: Ports are small and specific (\texttt{UserRepository}, \texttt{SessionRepository}).
  \item \textbf{Dependency Inversion}: The application layer depends on abstractions (ports), not storage details.
\end{itemize}

\subsection{DRY (Don't Repeat Yourself)}
Shared logic is centralized (e.g., \texttt{createCache} + \texttt{useApiCache}, reusable UI patterns such as modular settings fields in \texttt{Profile.tsx}).

\subsection{KISS (Keep It Simple, Stupid)}
The system uses straightforward state-based navigation and local persistence. This reduces operational complexity and makes the project easy to run and evaluate.

\subsection{Separation of Concerns}
\begin{itemize}[leftmargin=1.2em]
  \item UI rendering and interaction live in \texttt{src/components/}.
  \item Business rules live in \texttt{src/application/}.
  \item Storage details live in \texttt{src/infrastructure/}.
\end{itemize}

\subsection{Scalability \& Maintainability}
\begin{itemize}[leftmargin=1.2em]
  \item Clear layering allows adding a backend later (REST/GraphQL) while keeping domain and use cases stable.
  \item Repository pattern enables switching from browser storage to a database-backed persistence model.
  \item Typed domain models (\texttt{UserData}) reduce runtime errors and improve refactor safety.
\end{itemize}

% ------------------------------------------------------------
\section{Design Patterns Used}

\subsection{Dependency Injection (DI)}
Services are constructed once in \texttt{createServices()} and provided through React Context (\texttt{ServicesProvider}).
This separates construction from usage and makes it easy to swap implementations.

\subsection{Repository Pattern}
\texttt{UserRepository} and \texttt{SessionRepository} define the storage contract. Concrete repositories
(\texttt{LocalStorageUserRepository}, \texttt{SessionStorageSessionRepository}) implement browser persistence.

\subsection{Factory / Composition Root}
\texttt{createServices()} acts as a factory/composition root that assembles repositories and use cases.

\subsection{Observer / Reactive UI (React)}
React's state and hooks provide reactive updates. For example, dashboard views change based on local UI state and session storage values.

\subsection{MVC / MVVM (Relevance Note)}
The app is primarily a React SPA with component-driven UI rather than classic MVC. The closest mapping is:
\begin{itemize}[leftmargin=1.2em]
  \item \textbf{View}: React components
  \item \textbf{ViewModel/Controller}: Component state + handlers; use cases act as application controllers
  \item \textbf{Model}: \texttt{UserData} and stored session/history data
\end{itemize}

% ------------------------------------------------------------
\section{System Components Breakdown}

\subsection{Core Modules}
\begin{itemize}[leftmargin=1.2em]
  \item \textbf{\texttt{src/components/Dashboard.tsx}}: Main authenticated shell. Sidebar navigation and orchestration of feature views.
  \item \textbf{\texttt{src/components/InterviewSetup.tsx}}: Collects configuration (role, difficulty, type) and starts a session.
  \item \textbf{\texttt{src/components/InterviewRoom.tsx}}: Timed interview experience with webcam preview and simulated transcript generation.
  \item \textbf{\texttt{src/components/Report.tsx}}: Displays session report with metrics and comparison blocks.
  \item \textbf{\texttt{src/components/History.tsx}}: Lists past sessions from user history and renders progress indicators.
  \item \textbf{\texttt{src/components/PersonaLab.tsx}}: Persona selection UI (interviewer styles).
  \item \textbf{\texttt{src/components/ResumeAI.tsx}}: Resume analysis demo UI; gated behind file upload and analysis action.
  \item \textbf{\texttt{src/components/Profile.tsx}}: Account settings (modular fields, validation, save/reset actions).
  \item \textbf{\texttt{src/components/AboutUs.tsx}}: About page (also accessible from dashboard sidebar).
\end{itemize}

\subsection{Application Use Cases}
\begin{itemize}[leftmargin=1.2em]
  \item \textbf{Auth} (\texttt{src/application/useCases/auth.ts}): login, signup, get current user, logout.
  \item \textbf{Sessions} (\texttt{src/application/useCases/sessions.ts}): adds a session record, increments stats and skill scores.
  \item \textbf{User Profile} (\texttt{src/application/useCases/userProfile.ts}): updates user fields (name/password in demo).
\end{itemize}

\subsection{Infrastructure}
\begin{itemize}[leftmargin=1.2em]
  \item \textbf{Storage adapters} (\texttt{browserStorage.ts}): typed wrappers over localStorage/sessionStorage.
  \item \textbf{Repositories}: LocalStorage user repository and SessionStorage session repository.
\end{itemize}

% ------------------------------------------------------------
\section{Database Design}

\subsection{Type of Database Used}
This project uses \textbf{browser key-value storage} as persistence:
\begin{itemize}[leftmargin=1.2em]
  \item \textbf{localStorage} as a persistent user database.
  \item \textbf{sessionStorage} for session identity and navigation state.
\end{itemize}

\subsection{Schema Overview (Conceptual)}
The main persisted structure is a JSON object keyed by user email:
\begin{itemize}[leftmargin=1.2em]
  \item \textbf{User} (\texttt{UserData}): \texttt{email}, \texttt{name}, optional \texttt{password}, \texttt{stats}, \texttt{history}, \texttt{streak}, \texttt{skills}.
  \item \textbf{HistoryItem}: \texttt{id}, \texttt{role}, \texttt{date}, \texttt{score}, \texttt{type}.
\end{itemize}

\subsection{Relationships and Key Decisions}
\begin{itemize}[leftmargin=1.2em]
  \item A user has many history items (stored as an array inside the user record).
  \item Stats are embedded for fast reads on the dashboard (no joins in a key-value store).
  \item This design is suitable for demos and offline-first UX; a production system should migrate to a server database.
\end{itemize}

% ------------------------------------------------------------
\section{API Design}

\subsection{Current ``API'' (In-Process Service Calls)}
There is no HTTP API in the current version. The UI calls in-process use cases via a Services context:
\begin{itemize}[leftmargin=1.2em]
  \item \texttt{auth.signup(email, name)}
  \item \texttt{auth.loginByEmail(email)}
  \item \texttt{auth.getCurrentUser()}
  \item \texttt{sessions.addSession(email, historyItem, track)}
  \item \texttt{profile.updateUser(email, updates)}
\end{itemize}

\subsection{Authentication / Authorization Strategy}
\begin{itemize}[leftmargin=1.2em]
  \item \textbf{Current}: session identity is a stored email in sessionStorage.
  \item \textbf{Production-ready direction}: replace with token-based auth (JWT or opaque session id) managed by a backend, enforce authorization server-side, and avoid storing sensitive tokens in localStorage.
\end{itemize}

\subsection{Proposed REST Endpoints (Future Backend)}
For scalability, a future backend could expose endpoints such as:
\begin{itemize}[leftmargin=1.2em]
  \item \texttt{POST /auth/signup}, \texttt{POST /auth/login}, \texttt{POST /auth/logout}
  \item \texttt{GET /users/me}, \texttt{PATCH /users/me}
  \item \texttt{POST /sessions}, \texttt{GET /sessions}
  \item \texttt{POST /resume/analyze}
\end{itemize}

% ------------------------------------------------------------
\section{Scalability \& Performance Considerations}

\subsection{Load Handling}
The current app is client-only; performance mainly depends on browser rendering and local data size. For production:
\begin{itemize}[leftmargin=1.2em]
  \item Add backend services for AI generation, resume parsing, analytics computation.
  \item Use CDN for static assets; implement pagination for history.
\end{itemize}

\subsection{Caching}
The project includes a simple TTL cache (\texttt{src/lib/cache.ts}) and a hook wrapper (\texttt{src/hooks/useApiCache.ts}) suitable for caching API calls in a future backend version.

\subsection{Async Processing}
Current ``analysis'' tasks are simulated with timers. In a real system:
\begin{itemize}[leftmargin=1.2em]
  \item Resume analysis should run asynchronously (job queue) due to CPU/LLM latency.
  \item Websocket/SSE updates can provide real-time progress feedback.
\end{itemize}

% ------------------------------------------------------------
\section{Security Considerations}

\subsection{Authentication}
Current implementation uses sessionStorage for the logged-in email. This is adequate for a demo but not for production.

\subsection{Data Protection}
\begin{itemize}[leftmargin=1.2em]
  \item Avoid storing secrets (API keys) in the frontend. \texttt{geminiGenerateText} explicitly notes the need for a backend proxy.
  \item Treat user-provided inputs as untrusted; validate and sanitize on the server in a production version.
\end{itemize}

\subsection{Common Vulnerabilities (Notes)}
\begin{itemize}[leftmargin=1.2em]
  \item \textbf{XSS}: React escapes strings by default; still avoid injecting HTML from untrusted sources.
  \item \textbf{Sensitive storage}: do not store real passwords in localStorage for production; use hashed passwords server-side.
  \item \textbf{Permissions}: webcam access is user-approved; handle permission denial gracefully.
\end{itemize}

% ------------------------------------------------------------
\section{Diagrams Explanation}

\subsection{Class Diagram (\texttt{diagrams/Class\_diagram.png})}
\textbf{Purpose}: Illustrates the primary classes/modules and their relationships, especially around user/session data and core features.

\begin{figure}[h]
  \centering
  \includegraphics[width=\linewidth]{diagrams/Class_diagram.png}
  \caption{Class diagram showing core modules and relationships.}
\end{figure}

\subsection{Use Case Diagram (\texttt{diagrams/uml\_use\_case\_diagram.svg})}
\textbf{Purpose}: Captures functional interactions between the \textit{User} and the \textit{AI Mock Interview System}.

\textbf{Build note}: This diagram is stored as SVG. To compile this document with the diagram embedded, LaTeX typically requires \texttt{--shell-escape} and an SVG conversion tool (e.g., Inkscape). If your toolchain cannot render SVGs, export the SVG to PNG/PDF and update the include command.

\begin{figure}[h]
  \centering
  \includesvg[width=\linewidth]{diagrams/uml_use_case_diagram}
  \caption{Use case diagram: user and AI system interactions.}
\end{figure}

\subsection{Activity Diagram (\texttt{diagrams/uml\_activity\_diagram.png})}
\textbf{Purpose}: Shows the end-to-end control flow for interview setup, interview execution, and post-session reporting.

\begin{figure}[h]
  \centering
  \includegraphics[width=\linewidth]{diagrams/uml_activity_diagram.png}
  \caption{Activity flow for the interview process.}
\end{figure}

\subsection{Sequence Diagram (\texttt{diagrams/uml\_sequence\_diagram (1).png})}
\textbf{Purpose}: Visualizes time-ordered interactions (UI \(\rightarrow\) services \(\rightarrow\) repositories \(\rightarrow\) storage) across key user actions such as login and session completion.

\begin{figure}[h]
  \centering
  \includegraphics[width=\linewidth]{\detokenize{diagrams/uml_sequence_diagram (1).png}}
  \caption{Sequence diagram of core interactions.}
\end{figure}

\subsection{ER Diagram (\texttt{diagrams/Er\_Diagram.jpg})}
\textbf{Purpose}: Provides a conceptual data model. Although the current implementation is key-value storage, the ER diagram is a useful blueprint for migrating to a relational database.

\begin{figure}[h]
  \centering
  \includegraphics[width=\linewidth]{diagrams/Er_Diagram.jpg}
  \caption{Conceptual ER diagram for user/session-related data.}
\end{figure}

% ------------------------------------------------------------
\section{Challenges \& Trade-offs}

\subsection{Key Challenges}
\begin{itemize}[leftmargin=1.2em]
  \item Balancing a premium UI/UX with a maintainable component structure.
  \item Designing ``clean architecture'' boundaries in a frontend-only environment.
  \item Simulating AI features without exposing secrets or requiring a backend.
\end{itemize}

\subsection{Trade-offs}
\begin{itemize}[leftmargin=1.2em]
  \item \textbf{Client-only persistence}: simplifies setup and demos, but limits real security and multi-device access.
  \item \textbf{State-based navigation}: avoids routing complexity, but lacks deep-linking and URL shareability.
  \item \textbf{Mock AI}: good for UI prototyping; production requires secure server-side integration.
\end{itemize}

% ------------------------------------------------------------
\section{Future Improvements}
\begin{itemize}[leftmargin=1.2em]
  \item Add a backend (Node/Express or serverless) for real authentication, secure AI calls, and persistence.
  \item Replace localStorage ``DB'' with PostgreSQL (aligning with the ER diagram) and introduce migrations.
  \item Implement real STT via Web Speech API wrapper (\texttt{speechToTextOnce}) and store transcripts.
  \item Add role-based interview banks, question randomization, and more robust analytics.
  \item Introduce a router (React Router) for deep links (\texttt{/about}, \texttt{/dashboard}, etc.).
  \item Improve security posture: password hashing, session tokens, and CSP headers (server-side).
\end{itemize}

% ------------------------------------------------------------
\section{Conclusion}
AntriView AI demonstrates a portfolio-ready system design using a modern frontend stack, clean layering, and a maintainable approach to feature development.
While the current implementation is frontend-only, the architecture intentionally supports evolution into a full-stack system with secure AI services and persistent analytics.

\end{document}

